\pdfoutput=1
\documentclass[11pt]{article}
\usepackage[margin=0.98in]{geometry}
\usepackage{amsmath,amssymb,amsthm,mathtools}
\usepackage[hidelinks]{hyperref}
\hypersetup{pdftitle={Radiative Data Lost by the Newtonian Projection}}
\newtheorem{theorem}{Theorem}
\newtheorem{lemma}[theorem]{Lemma}
\newtheorem{proposition}[theorem]{Proposition}
\theoremstyle{remark}
\newtheorem{remark}[theorem]{Remark}
\newcommand{\dd}{\mathrm d}
\newcommand{\TT}{\mathrm{TT}}
\newcommand{\eps}{\varepsilon}
\newcommand{\tr}{\operatorname{tr}}
\newcommand{\supp}{\operatorname{supp}}
\newcommand{\Ric}{\operatorname{Ric}}
\newcommand{\Lie}{\mathcal L}
\newcommand{\R}{\mathbb R}
\newcommand{\Om}{\Omega}
\newcommand{\M}{m_{\mathrm{ADM}}}
\newcommand{\ME}{m_{\mathrm{esc}}}
\newcommand{\F}{\mathcal F}
\title{Radiative Data Lost by the Newtonian Projection\\
\large A finite-radius symplectic and ADM identity in Einstein--Euler}
\author{}
\date{}
\begin{document}
\maketitle
\begin{abstract}
The fast Newtonian limit of Einstein--Euler allows freely chosen initial gravitational waves at fixed limiting fluid data. We exhibit an exact constrained family with outgoing leading tangents in which the Newtonian projection forgets a rank-four radiative symplectic block and a positive mass Hessian. At a fixed vacuum extraction sphere, joint wave-amplitude derivatives of the actual metric converge in fast time after division by \(\varepsilon^2\), with \(\varepsilon=v_*/c\). The correspondingly normalized Einstein--Hilbert symplectic flux gives two nondegenerate canonical pairs; its time-translation energy pairing equals the leading wave-amplitude Hessian of the ADM mass of the same constraint solutions. The equality is at finite radius and uses no radiation-zone approximation or average. A delay-amplitude manifold gives an exact finite canonical sector of the resolved leading wave flow. The result is local in fast time and identifies a seed-dependent coefficient, not an exact nonlinear energy loss on a nonstationary tube, an asymptotic Bondi theorem for the fluid solutions, or a generic product splitting of the full theory.
\end{abstract}

\section{Setting and main result}

We use the asymptotically flat Einstein--Euler families in the regularized polytropic formulation of Oliynyk \cite{Oliynyk}. The material initial datum is one fixed smooth nonzero radial compact Makino density, with coordinate-rest velocity, supported inside \(r<a_0\). The term \emph{fixed material datum} refers to these coordinate Makino variables and the equation of state, not to a comparison at fixed proper baryon mass. All wave data below are smooth compact transverse trace-free (TT) tensor Cauchy data initially outside this matter, in a common shell \(a<r<b\) with \(a>a_0\). The exact metric and lapse/shift corrections from the relativistic constraints are not held fixed. The spatial coordinates and the fast coordinate \(s=t/\eps\) are those of the normalized harmonic construction; after the fixed unit change \(s\) is light-travel length and the limiting wave speed is one. \(\lambda\) denotes independent small real wave-amplitude coordinates, unrelated to \(\eps\).

We write formulas first in geometric normalization, with action coefficient \(1/(16\pi)\), so a mass has units of length. To recover physical mass multiply by \(c^2/G\), energy by \(c^4/G\), and a geometric symplectic form by \(c^3/G\). Spatial lengths are measured in a fixed physical unit; factors of that unit are understood in the integrals. Dots on \(H\) denote derivatives in fast time, not amplitude variation.

\begin{theorem}\label{thm:main}
There is a four-parameter local exact Einstein--Euler constraint and evolution family with the same slow Poisson--Euler limit and vacuum in one common marked finite extraction cylinder. At the zero-wave stock, for each parameter vector \(u\) the normalized first metric variation \(\eps^{-2}D_u g_\eps|_0\) tends in \(C^2\) on the fixed fast cylinder to a spatial TT outgoing solution \(H_u\), with error \(O(\eps)\). For the full outgoing aperture at a sufficiently large \emph{fixed} radius \(R\), the geometric Einstein--Hilbert symplectic flux on the exact tangents, divided by \(\eps^4\), tends to
\[
  \Omega(H_u,H_v)=\frac1{32\pi}\int_{\R^3}
                           (H_u:\dot H_v-H_v:\dot H_u)\,\dd^3x .
\]
The pullback on the four real tangents has rank four. The induced canonical-time energy pairing at the tube, normalized in the same way, is
\[
 m_{\mathrm{wave}}(H_u,H_v)=\frac1{64\pi}\int_{\R^3}
              (\nabla H_u:\nabla H_v+\dot H_u:\dot H_v)\,\dd^3x .
\]
It also equals
\(\lim_{\eps\to0}(2\eps^4)^{-1}D^2_\lambda\M|_0[u,v]\),
where \(\M\) is the geometric invariant ADM mass of the actual initial data. Both limits have error \(O(\eps)\) on fixed compact tangent sets. The equality is between these leading coefficients, not exact nonlinear conservation of a time-translation energy by the vacuum tube at finite \(\eps\).
\end{theorem}

\begin{proposition}\label{prop:delay}
At the resolved leading order there is a locally flow-invariant four-dimensional manifold of two independent nonzero real pulse amplitudes and center delays. Its symplectic form and positive wave Hamiltonian are
\[
      \Omega=\sum_{\sigma=+,\times}\dd(A_\sigma^2 m_\sigma)\wedge\dd\tau_\sigma,
           \qquad H_{\rm wave}=\sum_{\sigma=+,\times} A_\sigma^2m_\sigma,
\]
where \(m_\sigma>0\) is the wave mass of the corresponding unit pulse. This is a canonical completion of the projection at the fixed material preparation, not a claim of an exact nonlinear finite-dimensional truncation.
\end{proposition}

\paragraph{Position.}
Oliynyk already proves the fast Newtonian limit for a broad exact family and identifies its fast tensor radiation \cite[Prop.~4.1, Thm.~1.1]{Oliynyk}. We use that family, not a new derivation of the limit or an assertion that Newtonian source motion generates the outgoing pulse in this experiment. The new comparison is the nonzero \emph{bilinear} canonical data and ADM Hessian of an explicit outgoing seed family, obtained with simultaneous amplitude differentiation and an exact leading finite-radius flux identity. It differs from both a slow post-Newtonian binding-energy coefficient and the weak backreaction of a fixed-\(c\), small-wavelength metric family. Its news limit is taken only after extracting the coefficient at each finite radius.

\section{An exact constrained stock and its fast parameter jets}\label{sec:stock}

\subsection{Constraints with the material preparation held fixed}

Use the polytropic and Makino variables of \cite[Eqs.~(8)--(9), (29), (73)--(84)]{Oliynyk}; in particular a smooth compact Makino density allowed there produces smooth \(\rho,p\), and vacuum is an allowed value of the regularized fluid evolution. We need no physical-vacuum boundary theorem. Fix a nonzero radial \(\alpha_0\) in that class supported in \(r<a_0\) and fix \(w^I=0\) initially. Work with \(-1<\delta<-1/2\) and sufficiently high weighted Sobolev order \(k\); parameter seeds are smooth and compact, so the order can be chosen as high as required for finitely many trace derivatives.

The freely specified gravitational entries of the initial densitized inverse-metric variable \(U\) are \(z^{IJ}=\bar{\mathfrak u}^{IJ}\) and \(z_4^{IJ}=\partial_s\bar{\mathfrak u}^{IJ}\). Choose them from a small open finite-dimensional TT family, keeping the fluid entries fixed. The analytic constraint maps of \cite[Prop.~4.1]{Oliynyk} solve the Einstein, harmonic and four-velocity normalization constraints for the remaining entries. In their notation the leading \(44\) and mixed entries solve
\[
      \Delta\phi_0=\rho+\partial_I\partial_J z^{IJ},
      \qquad
      \Delta w^I_0=-\partial_J z_4^{JI}.
\]
Here \(w_0^I\) labels the mixed \emph{metric density} entry, not a change of the specified fluid velocity. Both extra terms vanish for our TT seeds. Thus \(\phi_0=\Delta^{-1}\rho\) is seed-independent and the decaying mixed entry is zero; the remaining initial \(\partial_s\) entries vanish at leading order by the harmonic constraints. These are genuine constrained data, not a linear wave added to an exact fluid metric without re-solving.

Write \(g_\eps(\lambda)\) for the physical metric in the fast coordinates, denoting ordinary spatial components by \(I,J\). Its reconstruction from \(U\), in the conventions with \(\eta=\operatorname{diag}(1,1,1,-1)\) and \(s\) last, begins with
\[
 g_{\eps,ij}=\eta_{ij}
  +2\eps^2\left(\eta_{k\ell}U^{k\ell}\eta_{ij}
                     -2\eta_{ik}U^{k\ell}\eta_{\ell j}\right)
  +O(\eps^4)
\]
when \(U\) is bounded. The \(O(\eps)\) errors in the solution's \(U\) turn into \(O(\eps^3)\) here. Thus if \(u,v\) are any fixed amplitude directions at zero wave and \(H_u\) has compact TT initial seeds \(-4D_u(z,z_4)\), then at \(s=0\),
\begin{align}
   D_u g_\eps&=\eps^2 H_u+O(\eps^3),&
   D_u \partial_s g_\eps&=\eps^2\dot H_u+O(\eps^3),\label{eq:initfirst}\\
   D_uD_v g_\eps&=O(\eps^3),&
   D_uD_v\partial_s g_\eps&=O(\eps^3).\label{eq:initsecond}
\end{align}
The leading tensors have only spatial components. The estimates hold in \(H^k_\delta,H^{k-1}_{\delta-1}\) for the first and time-derivative entries as appropriate, and after equivalent spatial differentiation. The initial shift is \(O(\eps^3)\) with two amplitude derivatives. For the linear coordinates at zero wave used in the theorem, this is a joint Taylor statement from the analytic dependence on the initial seed entries, since the \(\eps^2\) metric coefficient is affine in those entries. It is not a derivative of \cite[Thm.~1.1]{Oliynyk}'s value-level slow-limit estimate. For the delay family used later the constraints are evaluated on a smooth translated seed path; the composition is smooth, not generally analytic in the delays.

\subsection{Fast evolution and simultaneous parameter estimates}

Use the original \emph{local} reduced Einstein--Euler equation \cite[Eq.~(64)]{Oliynyk}, which the exact solution satisfies while the compact fluid remains inside the limiting support bound. With \(s=t/\eps\), write its variables schematically as
\(V_\eps=(U_4,U_I,\eps(U-U(0)),\alpha,w)\), with \(U_4=\partial_s U\), \(U_I=\partial_I U\). Multiplying the original equation by \(\eps\) gives the symmetric form
\begin{equation}
\begin{aligned}
 b^0(\eps V_\eps,\eps^2U_\eps(0))\partial_s V_\eps
 &=c^I\partial_I V_\eps+\mathcal G(V_\eps)
     +\eps\,b^I_\eps(V_\eps,U_\eps)\partial_I V_\eps
     +\eps\mathcal F_\eps(V_\eps,U_\eps),\\
 U_\eps(s)&=U_\eps(0)+\int_0^s U_{4,\eps}(s')\,\dd s' .
\end{aligned}
 \label{eq:fasteq}
\end{equation}
In the \(b^0\) argument \(U_\eps(0)\) occupies the initial-density slot of Oliynyk's block vector: together with \(\eps V_\eps\) it reconstructs \(\eps^2U_\eps(s)\), without counting the increment twice. The temporal \(b^0\) is uniformly positive on a fixed neighborhood of bounded fast data and tends to the identity in these variables. The \(c^I\) are the constant symmetric principal wave matrices on the \(U_4,U_I\) block, with zero entries in the frozen blocks. The only order-one source \(\mathcal G\) has the component \(-\rho(\alpha)\) in the \(44\) temporal wave row; it depends on no amplitude except through \(\alpha\). The schematic \(b^I_\eps,\mathcal F_\eps\) collect all remaining regular arguments using the same reconstruction. They are regular including across \(\alpha=0\), with \(\eps\)-weighted lower terms as in \cite[Eqs.~(65)--(72)]{Oliynyk}. In particular the \(\eps=0\) fluid block is frozen. The unscaled metric variable must be reconstructed by the second formula of \eqref{eq:fasteq}; one cannot read its variation from the frozen scaled increment alone.

At \(\eps=0\) the fixed source and its fixed initial \(\phi_0\) give a stationary \(44\) block because \(\Delta\phi_0=\rho\). The first \(\lambda\)-derivative of the fluid block is zero, so all first amplitude derivatives of the metric density are free tensor waves with the given initial data. Their divergence, trace and time components stay zero by the flat wave equation and the TT initial conditions. In the linear seed coordinates all second amplitude derivatives of the limiting system are zero. This is the same fast branch, now linearized in the independent seed variables.

For clarity the parameter comparison can be shown directly by fixed-slab energies. Choose any fixed compact fast slab \(0\leq s\leq S\) and a ball containing the eventual extraction tube and the domains of dependence of its collars, with a fixed margin for the bounded characteristic speed. Smooth data permit standard commuted energies on this enlarged region (or energies with a moving outer domain of dependence). In a neighborhood of the bounded \(\eps=0\) solution, \eqref{eq:fasteq} has a uniform symmetrizer and finite coefficient bounds; the density reconstruction is controlled by its displayed time integral. At one lower Sobolev order the difference between the \(\eps\)-solution and the \(\eps=0\) solution has initial \(O(\eps)\) norm and \(O(\eps)\) inhomogeneity from the coefficient changes applied to the smooth limiting solution. Sobolev multiplication and Gr\"onwall on the fixed interval then give \(O(\eps)\) differences and provide continuation in that neighborhood. The exact existence interval in slow time is available from \cite[Prop.~4.3, Thm.~1.1]{Oliynyk}, hence contains the fast slab for sufficiently small \(\eps\).

Differentiating \eqref{eq:fasteq} in a seed coordinate gives a linear symmetric hyperbolic principal part, with lower terms involving the smooth unvaried solution and the reconstructing time integral. The second seed derivative has the same principal part and a source bilinear in first derivatives and their first spatial derivatives. The ordinary difference-quotient argument and commuted energy bound give these derivatives in successively lower Sobolev orders. Subtract their \(\eps=0\) equations: initial differences are \(O(\eps)\) by the joint constraint Taylor expansion, and the forces are \(O(\eps)\) at the next lower order by the already bounded lower jets and regular coefficient differences. Taking sufficiently many initial derivatives before these finitely many losses, Sobolev embedding and \eqref{eq:fasteq} for the time jets give on the chosen fixed tube and caps
\begin{equation}
             D_u g_\eps=\eps^2 H_u+O_{C^2}(\eps^3).
             \label{eq:fasttangent}
\end{equation}
The estimates are joint on compact small seed sets and uniform on bounded amplitude directions. They do not run an energy estimate to \(s\) of order \(\eps^{-1}\). The Makino equation is transport of \(\alpha\) along the fluid, with a source proportional to \(\alpha\) itself. In fast coordinates its support velocity is \(\eps w^I/(1+\eps w^4)=O(\eps)\), so the density is supported inside \(a_0+O(\eps)\) on the fixed interval. This is the support step of \cite[Prop.~4.3]{Oliynyk}. The auxiliary spatial cutoffs in that construction may be chosen equal to one throughout the fixed observation region and its domain of dependence. Thus a fixed extraction cylinder separated from the original material is actually vacuum, with a common open collar for all sufficiently small seed parameters.

The slow Poisson--Euler limit with these initial variables is the same for every seed by the fluid convergence and uniqueness in \cite[Thm.~1.1]{Oliynyk}. The fast comparison just proved uses that exact family but yields a different observation: the normalized first tensor derivative on the fast tube is nonzero despite the common material limit.


\section{Outgoing waves and the finite extraction identity}\label{sec:radiation}

\subsection{Compact outgoing TT waves in both polarizations}

Choose any nonzero \(f\in C^\infty_c((a,b))\). For \(s\geq0\) set
\(\phi(s,x)=f(r-s)/r\) for \(r=|x|>0\), extended by zero near the origin. It is an exact smooth solution of the scalar flat wave equation since its numerator obeys the one-dimensional radial wave equation and its support stays away from \(r=0\). Let
\[
 P_{IJ}=\delta_{IJ}\Delta-\partial_I\partial_J,\qquad
 (T_A)_{IJ}=\left(P_{IK}P_{JL}-\tfrac12 P_{IJ}P_{KL}\right)A_{KL}
\]
for a constant symmetric matrix \(A\). For \(H_A=T_A\phi\) the divergence vanishes by \(\partial_I P_{IJ}=0\), and the trace by \(P^2=\Delta P\), \(\tr P=2\Delta\). It solves the flat tensor wave equation because the operators are constant coefficient. Its exact support is contained in \(a+s<r<b+s\). In particular the pure outgoing property here is not an assumption based on large radius or on a high-frequency expansion.

Use \(A_+=\operatorname{diag}(1,-1,0)\) and \(A_\times\) with \(xy=yx=1\) and all other entries zero. In the radiation field at direction \(e_z\), their limiting screen amplitudes \(rH_A(r+u,rn)\) are \(f^{(4)}(-u)A_+\) and \(f^{(4)}(-u)A_\times\): each leading radial derivative is accompanied by the corresponding direction cosine, and \(P\) projects into the transverse plane. The fourth derivative is not identically zero. Thus these supply two distinct real polarizations, whose complex combinations also give the circular helicities. Write \(H_\sigma=H_{A_\sigma}\). For each \(\sigma\) use the two seed waves \(H_\sigma,\dot H_\sigma\) in the four real parameter stock. Both are compact TT with their time derivatives initially compact. By \eqref{eq:initfirst} their exact families are obtained by the corresponding \(-1/4\) rescalings of the densitized seeds. These are actual independently varied real wave data, not one solution with four choices of normalization.

The entire family is invariant under spatial inversion as a tensor field: \(H_\sigma(-x)=H_\sigma(x)\) in its spatial rank-two components, and the time derivatives have the same evenness. The matter preparation and \(\lambda\)-family constraint equations share this symmetry. The analytic constraint solution is the local unique branch, so the exact initial metric and extrinsic curvature have tensorial inversion symmetry (an initial mixed spatial/temporal component changes sign). Consequently the exact total ADM spatial momentum is zero at all sufficiently small \(\lambda\); the ADM energy in geometric normalization equals the positive invariant mass. No leading calculation of the ADM momentum is used to impose this conclusion.

\subsection{Einstein--Hilbert normalization on a finite vacuum tube}

Fix \(R>b\) with a positive margin, and a common fast-time aperture containing the crossing times \([R-b,R-a]\) with margins. Its early and late sections can be chosen at positive times before and after the crossing. At the early section the complete leading waves lie inside the sphere, at the late section they lie outside, and a time neighborhood of each end has no leading wave on the aperture. Use the full sphere rather than one angular patch. The exact background at \(\lambda=0\) is \(C^2\)-close to Minkowski by \(O(\eps^2)\) on this fixed cylinder. The radius of the marked coordinate sphere can be adjusted by \(O(\eps^2)\) on that background to have area \(4\pi R^2\); the area derivative at the flat sphere is \(8\pi R\), and the regular fast derivatives make the induced worldtube and conormal change of this same order. We evaluate on this marked worldtube, using the exact tangent fields pulled back to it. This calibration, or a fixed coordinate sphere with the same limit, changes none of the leading current coefficients. The whole tube is vacuum by Section~\ref{sec:stock}.

For completeness the quadratic Einstein--Hilbert TT action at Minkowski, in the geometric convention, is
\begin{equation}
 \frac1{64\pi}\int\bigl(|\dot H|^2-|\nabla H|^2\bigr)\,\dd s\,\dd^3x .
 \label{eq:ttaction}
\end{equation}
Indeed the linearized Einstein tensor on spatial TT is \(-\Box H/2\), and the second action coefficient before integration by parts is \(-\int H:G^{(1)}(H)/(32\pi)\), where \(\Box=-\partial_s^2+\Delta\). Thus the symplectic current from \eqref{eq:ttaction} is the restriction of the Einstein--Hilbert current, up to the usual exact divergence from the variational potential. We choose its sign so its Cauchy pullback is \(\int\delta H\wedge\delta(\dot H/(32\pi))\), and write outward flux with this convention. In a smooth potential representative such as the Einstein--Hilbert variational potential of \cite{IyerWald}, a representative change integrates to a boundary term on this tube. Its leading term vanishes at the two caps since the leading wave fields and their derivatives have support away from them, and there are no angular edges. This justifies using the simple TT expression even though the extraction radius is finite.

The resulting flat currents give
\begin{align}
 \Omega(H_1,H_2)
   &= \frac1{32\pi}\int_{\R^3}
        (H_1:\dot H_2-H_2:\dot H_1)\,\dd^3x, \label{eq:Omega}\\
 \F_{\Omega,R}(H_1,H_2)
   &= -\frac1{32\pi}\int_{\rm aperture}R^2
        (H_1:\partial_r H_2-H_2:\partial_r H_1)\,\dd s\,\dd\Omega, \label{eq:symout}\\
 \ME(H_1,H_2)
   &=\frac1{64\pi}\int_{\R^3}
      (\dot H_1:\dot H_2+\nabla H_1:\nabla H_2)\,\dd^3x .
       \label{eq:masswave}
\end{align}
Here \(\dd\Omega\) is the unit sphere measure and the spatial integrals may be taken on any fixed fast slice. The time component and spatial flux density of the symplectic current satisfy continuity by the wave equation. Thus on the outgoing wave stock \(\F_{\Omega,R}=\Omega\): the interior spatial form changes from the full value to zero as the pulse crosses. For the energy the local flat continuity equation has outward density
\(-(\dot H_1:\partial_rH_2+\dot H_2:\partial_rH_1)/(64\pi)\).
It follows, again by the two cap supports, that
\begin{equation}
 \ME(H_1,H_2)
  =-\frac1{64\pi}\int_{\rm aperture} R^2
       (\dot H_1:\partial_r H_2+\dot H_2:\partial_r H_1)
                      \,\dd s\,\dd\Omega .
 \label{eq:energyout}
\end{equation}
The sign is outward and positive on the fully escaped wave when \(H_1=H_2\). This identity uses the exact finite-radius fields, including their radial derivative, not the far-zone replacement \(\partial_r=-\partial_s\).

To state the comparison at finite \(\eps\) without inventing a conserved nonlinear energy, let \(h_{\eps,u}=D_u g_\eps|_0\) and let \(\F^\eps_{\rm EH}\) be the full signed geometric Einstein--Hilbert symplectic current integrated over the above vacuum worldtube at the zero-wave solution. For wave parameters define the canonical-time energy diagnostic
\[
 \mathcal E_\eps(u,v)
  =-\tfrac14\left[
    \F^\eps_{\rm EH}(h_{\eps,u},\Lie_{\partial_s}h_{\eps,v})
   +\F^\eps_{\rm EH}(h_{\eps,v},\Lie_{\partial_s}h_{\eps,u})\right] .
\]
The Lie derivative in this expression is just the chosen time reading; at finite \(\eps\) the background need not be stationary and the second argument need not be a linearized solution. This is a well-defined local bilinear current diagnostic, not a claim of nonlinear balance under a missing Killing field. The smooth current is first order in its arguments and bilinear. Equation~\eqref{eq:fasttangent}, with the \(C^2\) error for time-translation arguments, gives
\begin{align}
 \F^\eps_{\rm EH}(h_{\eps,u},h_{\eps,v})
      &=\eps^4\Omega(H_u,H_v)+O(\eps^5),\label{eq:EHlimit}\\
 \mathcal E_\eps(u,v)&=\eps^4\ME(H_u,H_v)+O(\eps^5).
      \label{eq:Elimit}
\end{align}
For the second line one integrates \eqref{eq:symout} by parts in \(s\), with no leading cap term; \(-\tfrac14\) times its two time-translation pairings is precisely \eqref{eq:energyout}. Perturbations of the background coefficients and tube geometry are higher order. This uses an actual quadratic diagnostic of the \emph{same} exact Einstein--Euler solutions, not an identification of the unscaled gravitational energy with a fluid observable.

The symplectic block is nondegenerate: reflection \(y\mapsto -y\) fixes \(H_+\) and negates \(H_\times\) as tensors, so all cross-polarization energy and symplectic pairings vanish (also at different fixed delays). For either \(\sigma\) the wave equation gives
\[
      \Omega(H_\sigma,\dot H_\sigma)
       =-\frac1{32\pi}\int_{\R^3}
                  (|\nabla H_\sigma|^2+|\dot H_\sigma|^2)\,\dd^3x
       =-2m_\sigma,\qquad m_\sigma=\ME(H_\sigma,H_\sigma)>0 .
\]
The two real quadrature pairs thus give rank four. The flat linearized curvature component is \(C^{(1)}_{sIsJ}=-\partial_s^2H_{IJ}/2\), which is gauge invariant on the flat background. If it vanished on the full aperture for a nonzero tangent, the initial zero margin would give \(\partial_sH=0\) there throughout. This contradicts the strictly positive outgoing energy \eqref{eq:energyout}. Thus the retained directions are not pure gauge even on the finite extraction tube.

\subsection{A finite canonical amplitude--delay manifold}

The fixed linear four-dimensional tangent span need not stay invariant under the whole fast evolution. Instead take \(H(s)=\sum_{\sigma}A_\sigma H_\sigma(s-\tau_\sigma)\) with the delays restricted to a small interval and \(A_\sigma\ne0\). The original scalar pulses extend smoothly to this small backward interval because their support still stays away from the origin; enlarge the initial shell and extraction window by a fixed small margin. These are admissible smooth seed paths of the same exact Einstein--Euler constraint map. Reflection eliminates mixed-polarization pairings for all delays. Since \(\partial_{\tau_\sigma}H=-A_\sigma\dot H_\sigma\) and \(\partial_{A_\sigma}H=H_\sigma\), the symplectic form is \(2A_\sigma m_\sigma\,\dd A_\sigma\wedge\dd\tau_\sigma\) on each pair. The energy is \(A_\sigma^2m_\sigma\). Thus on every nonzero-amplitude chamber this form is nondegenerate and the Hamiltonian vector field translates \(\tau_\sigma\) to \(\tau_\sigma-s\) when \(\iota_X\Omega=\dd H_{\rm wave}\). It is exactly the local leading wave evolution on this four-dimensional manifold. At a zero pulse there is no delay coordinate; the previous two linear quadratures per polarization give the nondegenerate tangent result through zero. This proves Proposition~\ref{prop:delay}.


\section{The ADM Hessian of the actual initial family}\label{sec:adm}

\subsection{The fixed-source terms do not enter this Hessian}

Use \(\gamma_{IJ}=\delta_{IJ}+h_{IJ}\) for the exact initial spatial metric, and \(K\) for its extrinsic curvature. With the unit normal \(n\) and lapse tending to one at infinity, the exact Hamiltonian constraint in geometric normalization is
\begin{equation}
          R(\gamma)+(\tr_\gamma K)^2-K:_{\gamma} K=16\pi\mu_{\rm geom}.
          \label{eq:HC}
\end{equation}
The source \(\mu_{\rm geom}\) here is normal energy, not the coordinate component of the four-velocity. In the dimensionless normalization \(G^{ij}=2\eps^4 T^{ij}\) of \cite[Eq.~(1)]{Oliynyk}, put \(u=\eps v\) for the unit fluid four-velocity after the fast coordinate change. At \(s=0\) the coordinate-rest condition makes \(u\) proportional to \(\partial_s=Nn+\beta\), so
\[
        W^2=(u\cdot n)^2=\frac{1}{1-|\beta|_\gamma^2/N^2},\qquad
        \mu_{\rm geom}
          =\frac{\eps^2}{4\pi}
                \bigl[\rho+(\rho+\eps^2p)(W^2-1)\bigr].
\]
The second formula follows by \(16\pi\mu_{\rm geom}=4\eps^4 T_{nn}\) and the velocity normalization in the stress tensor. On the fixed compact material support, \(\beta=O(\eps^3)\) with two \(\lambda\)-derivatives, from the exact joint initial construction; \(N=1+O(\eps^2)\). Thus
\begin{equation}
           \mu_{\rm geom}=\frac{\eps^2\rho}{4\pi}+O(\eps^8)
           \label{eq:sourceexp}
\end{equation}
jointly with up to two amplitude derivatives of the remainder. The leading \(\rho,p\) in these coordinates are prescribed exactly. Equation~\eqref{eq:sourceexp} does not assume the fluid is exactly normal after the relativistic shift is solved. It is also not a statement at fixed proper fluid volume or proper baryon number, either of which can vary through the exact \(\gamma\).

\subsection{Differentiation by an integrated constraint identity}

Let \(Lh=\partial_I\partial_Jh_{IJ}-\Delta\tr h\), with Euclidean contractions. The standard ADM mass in the inversion-symmetric family is
\[
 \M=\frac1{16\pi}\lim_{\ell\to\infty}
        \int_{S_\ell}(\partial_J h_{IJ}-\partial_I\tr h)n^I\,\dd S .
\]
Use \eqref{eq:HC} to evaluate \(\int Lh\) and then differentiate on the parameter family. For the curvature write
\[
       R(\delta+h)=Lh+A(h)\partial^2h+B(h)(\partial h)^2
\]
schematically, with tensor contractions and smooth analytic coefficients and \(A(0)=0\). After integration by parts in the second term, the nonlinear curvature integral \(Q(h)=\int(R-Lh)\,\dd^3x\) has only products of first derivatives with smooth metric coefficients. More precisely the exterior products \(h\partial h\) on \(S_\ell\) vanish as \(\ell\to\infty\), since the weighted Sobolev estimates yield \(h=O(r^\delta)\), \(\partial h=O(r^{\delta-1})\) and \(2\delta+1<0\). Both \(\partial h\) and \(K\) are in \(L^2(\dd^3x)\) in these weighted orders. Thus the integrated quadratic and higher terms exist, and the limit of the integrated linear curvature is fixed by the exact scalar constraint. We do not require \(Lh\) to be \(L^1\) for arbitrary metrics in the weighted space. In particular for the exact family,
\begin{equation}
 16\pi\M
    =\int_{\R^3}\bigl(16\pi\mu_{\rm geom}
                 +K:_{\gamma}K-(\tr_\gamma K)^2\bigr)\,\dd^3x-Q(h).
 \label{eq:intadm}
\end{equation}
The Euclidean volume in this formula is intentional, as required by the flat ADM divergence. The right side has smooth joint amplitude derivatives by the compact source support and the weighted product estimates for the integrated terms. This avoids assuming continuity of the ADM surface expression on any larger arbitrary weak topology.

For a compact TT \(H\), the quadratic part of \(Q\) is
\begin{equation}
          Q_2(H)=-\tfrac14\int_{\R^3}|\nabla H|^2\,\dd^3x .
          \label{eq:Q2}
\end{equation}
Indeed \(\int\sqrt{\gamma}R\) on a covariant metric variation has quadratic coefficient \(-\tfrac12\int H:G^{(1)}(H)\), where \(G^{(1)}=-\Delta H/2\) in spatial TT gauge. Integrating by parts yields \eqref{eq:Q2}. The determinant correction at quadratic order pairs with the linear scalar term; this term vanishes on TT. Thus it is the same coefficient for \(Q\) with Euclidean measure. Its mixed version follows by polarization.

On the exact family \(h,K=O(\eps^2)\) initially. With the extrinsic curvature sign \(K=-\tfrac12\Lie_n\gamma\), the expansions \eqref{eq:initfirst}--\eqref{eq:initsecond} and initial shift estimate imply
\[
       D_u K=-\tfrac12\eps^2\dot H_u(0)+O(\eps^3),\qquad
       D_uD_vK=O(\eps^3).
\]
Only the absolute coefficient matters below. The leading trace of the first variation is zero. Every term in the Hessian of the curvature and extrinsic nonlinearities with one background factor and one second amplitude derivative is \(O(\eps^5)\). Terms from a higher derivative order in the metric amplitude or index-raising corrections are also higher order. This follows directly from the joint data estimates and the integrable two-first-derivative representation of \(Q\). Differentiating \eqref{eq:intadm}, dividing the Hessian by two, and using \eqref{eq:sourceexp} and \eqref{eq:Q2} gives
\begin{equation}
 \frac12D^2_\lambda\M|_0[u,v]
   =\frac{\eps^4}{64\pi}\int_{\R^3}
       (\nabla H_u(0):\nabla H_v(0)+\dot H_u(0):\dot H_v(0))\,\dd^3x
       +O(\eps^5).
 \label{eq:ADMhess}
\end{equation}
It equals the leading outgoing finite-tube coefficient in \eqref{eq:Elimit}. Every retained term came from the actual exact constraints; \eqref{eq:ADMhess} does not follow by differentiating the published \(O(\eps^2)\) estimate for the slow-unit ADM mass. It isolates the coefficient quadratic in freely set wave data, not the entire binding, pressure, or other correction of the slow theory. This completes the proof of Theorem~\ref{thm:main}.

\section{Units, the news check, and limitations}

For a single real amplitude \(\lambda\), the escaped mass coefficient \([\lambda^2]\) is
\((c^2/G)\eps^4\ME(H,H)\) at the retained order, and the energy coefficient is \((c^4/G)\eps^4\ME(H,H)\). The symplectic normalization is \((c^3/G)\eps^4\). At fixed characteristic fluid speed \(v_*\), \(c=v_*/\eps\); the unscaled geometric coefficient vanishes, physical mass is order \(\eps^2\), and physical symplectic form is order \(\eps\). The physical energy coefficient can be order one; it is still invisible to the Poisson--Euler projection of the same fixed material data. Combining it consistently with the symplectic action factor uses physical time \(s/c\), not the same numerical normalizing factor for the two physical flows.

The standard news check can be obtained after the finite-radius coefficient extraction. For fixed \(u=s-r\) let \(C(u,n)=\lim_{r\to\infty}rH(r+u,rn)\); it is the transverse trace-free radiation field of our explicit compact outgoing profile. Put \(N=\partial_u C\). Then \(\partial_sH=r^{-1}N+O(r^{-2})\) and \(\partial_rH|_s=-r^{-1}N+O(r^{-2})\) in the supported outgoing region. Hence the diagonal \eqref{eq:energyout} tends to \((32\pi)^{-1}\int N:N\,\dd u\,\dd\Omega\), the positive escape convention, the negative of the final-minus-initial Bondi mass change in that normalization \cite{Bondi}. This is an iterated limit of the extracted flat coefficient. No estimate uniform in \(R\) as \(\eps\to0\), asymptotic completeness of the fluid solutions, or nonlinear Bondi statement for them is used.

The wave amplitudes are independent initial radiation, not the far-zone quadrupole of the slow matter. Rigorous weak-field/slow-motion radiation with matter-sourced asymptotics, for example \cite{FutamaseSchutz}, addresses a different preparation and scale. Here the exact finite sphere sees a rank-four form and an ADM Hessian lost by the fixed-material Newtonian projection. The construction neither supplies a global nonlinear slow/fast symplectic splitting with variable fluids nor claims the existence of an instantaneous Newtonian observable reproducing the outgoing waveform.


\begin{thebibliography}{9}
\bibitem{Oliynyk}
T.~A.~Oliynyk,
The fast Newtonian limit for perfect fluids,
\emph{Adv. Theor. Math. Phys.} \textbf{16} (2012), 359--391,
\href{https://arxiv.org/abs/0908.4455}{arXiv:0908.4455}.
\bibitem{IyerWald}
V.~Iyer and R.~M.~Wald,
Some properties of Noether charge and a proposal for dynamical black hole entropy,
\emph{Phys. Rev. D} \textbf{50} (1994), 846--864,
\href{https://arxiv.org/abs/gr-qc/9403028}{arXiv:gr-qc/9403028}.
\bibitem{Bondi}
H.~Bondi, M.~G.~J.~van der Burg and A.~W.~K.~Metzner,
Gravitational waves in general relativity. VII. Waves from axi-symmetric isolated systems,
\emph{Proc. R. Soc. Lond. A} \textbf{269} (1962), 21--52,
\href{https://doi.org/10.1098/rspa.1962.0161}{DOI:10.1098/rspa.1962.0161}.
\bibitem{FutamaseSchutz}
T.~Futamase and B.~F.~Schutz,
Gravitational radiation and the validity of the far-zone quadrupole formula in the Newtonian limit of general relativity,
\emph{Phys. Rev. D} \textbf{32} (1985), 2557--2565,
\href{https://journals.aps.org/prd/abstract/10.1103/PhysRevD.32.2557}{DOI:10.1103/PhysRevD.32.2557}.
\end{thebibliography}


\end{document}

